\subsection{MATRICES OVER A FIELD}

The finite matrices with m rows and n columns over a field K are in one-to-one correspondence with the linear mappings of the right vector space \(E = K_{d}^{n}\) into the right vector space \(K_{d}^{m}\) when the matrices of these mappings are taken with respect to the canonical bases of E and F. By definition, the rank of such a matrix X is the rank of the linear mapping \(u: E \to F\) corresponding to it; as this number is by definition the dimension of the subspace \(u(E)\) of F, it amounts to the same (identifying the columns of X with the images under u of the canonical basis of E) to give the following definition:

DEFINITION 7. Given a matrix X with m rows and n columns over a field K, the dimension of the subspace of \(K_{d}^{m}\) generated by the n columns of X is called the rank of X with respect to K and denoted by \(\operatorname{rg}(X)\) .

It can also be said that the rank of X is the maximum number of linearly independent columns of X (as elements of \(K_{d}^{m}\) ). Obviously \(\operatorname{rg}(X) \leqslant \inf(m, n)\) ; for every submatrix Y of X, \(\operatorname{rg}(Y) \leqslant \operatorname{rg}(X)\) .

If E and F are two finite-dimensional vector spaces over K and u a linear mapping of E into F, the rank of the matrix \(M(u)\) with respect to any two bases is equal to the rank of u.

PROPOSITION 9. If the elements of a matrix X with m rows and n columns belong to a subfield \(K_{0}\) of a field K, the rank of X with respect to \(K_{0}\) is equal to the rank of X with respect to K.

Let \(F_{0}\) be the right vector \(K_{0}\) -space generated by the canonical basis of the right vector K-space \(E = K_{d}^{m}\) ; by hypothesis the columns of X belong to \(E_{0}\) . Let \(V_{0}\) (resp. V) be the vector sub- \(K_{0}\) -space of \(F_{0}\) (resp. the vector sub-K-space of E) generated by these columns. Then \(V = V_{0} \otimes_{K_{0}} K\) (§ 8, no. 2, Proposition 2) and hence \(\dim_{K} V = \dim_{K_{0}} V_{0}\) .

PROPOSITION 10. The rank of a matrix X over a field K is equal to the rank of its transpose \(^{t}X\) over the opposite field \(K^{0}\) .

In the notation introduced before Definition 7, the rank of u is equal to that of \({}^{t}u\) (§ 7, no. 5, Proposition 10) and the proposition therefore follows from no. 4, Proposition 3.

It is thus seen that the rank of X can also be defined as the maximum number of linearly independent rows of X (considering them as elements of the left vector K-space \(K_{s}^{n}\) ).

The square matrices of order n over a field K correspond bijectively with the endomorphisms of \(E = K_{d}^{n}\) and form a ring isomorphic to the ring

End \(_{\kappa}\) (E) (no. 7); corresponding to the automorphisms of E are the invertible square matrices.

PROPOSITION 11. Let \(X\) be a square matrix of order \(n\) over a field \(K\) . The following properties are equivalent:

\begin{enumerate}
\def\labelenumi{(\alph{enumi})}
\item
  \(X\) is invertible in \(\mathbf{M}_n(\mathbf{K})\) .
\item
  \(X\) is right invertible in \(\mathbf{M}_n(\mathbf{K})\) .
\item
  \(X\) is left invertible in \(\mathbf{M}_n(\mathbf{K})\) .
\item
  \(X\) is of rank \(n\) .
\end{enumerate}

This is just a translation of § 7, no. 4, Corollary to Proposition 9.

PROPOSITION 12. For a system of \(m\) linear equations in \(n\) unknowns

\[
\sum_ {j = 1} ^ {n} a _ {i j} x _ {j} = b _ {i} \quad (1 \leqslant i \leqslant m)\tag{54}
\]

over a field K to have at least one solution, it is necessary and sufficient that the matrix \(A = (a_{ij})\) of the system and the matrix B, obtained by bordering A with an \((n + 1)\) -th column equal to \((b_{i})\) , be matrices of the same rank.

It has been seen (no. 4) that the existence of a solution of (54) is equivalent to the fact that the column \((b_{i})\) is a linear combination of the columns of A and the proposition therefore follows from §7, no. 3, Corollary 4 to Proposition 4.

Note that the condition of Proposition 12 is always fulfilled when m = n and A is invertible, that is of rank n (Proposition 11). If x and b then denote the matrices with one column \((x_{i})\) and \((b_{i})\) respectively, system (54) is equivalent to A. x = b and its unique solution is \(x = A^{-1}.b\) .

